\section*{附录}

\subsection*{A. 软件环境与复现说明}

本文采用Python 3.12.13完成计算，使用NumPy、pandas、Matplotlib、Seaborn、OpenPyXL、SciPy和HiGHS完成科学计算、数据处理、绘图和优化求解，使用\texttt{uv}管理运行环境。依赖约束和锁定信息见支撑材料中的\texttt{pyproject.toml}与\texttt{uv.lock}。

复核时先同步运行环境，再按问题一至问题四的顺序执行题目级入口。问题二至问题四的预处理依赖问题一生成的共享处理层，因此不将项目根目录的\texttt{main.py}作为无条件复现入口。运行命令如下：

\begin{lstlisting}[language={},basicstyle=\ttfamily\small]
uv sync
uv run python question\question_01\main.py
uv run python question\question_02\main.py
uv run python question\question_03\main.py
uv run python question\question_04\main.py
\end{lstlisting}

各题入口依次完成数据处理、模型求解、模型检验和结果绘图。支撑材料中同时保留了处理后的模型输入和已生成的结果表。若仅复核论文中的数值，可直接检查结果表及独立检验表；若要从原始附件重新计算，应在保持题目数据不变的前提下重新完成数据处理。

\subsection*{B. 源程序与支撑材料}

支撑材料包含本题全部计算机源程序、运行环境配置、处理后的模型输入、问题一至问题四的关键结果表、模型检验表、论文图形以及文件校验清单。结果表用于追溯正文中的具体数值，图形用于论文展示，独立检验表用于核对任务覆盖、资源容量、能源平衡、储能状态和时间边界等条件。

各问题的支撑结果分别包括：问题一的需求预测、基础算力调度和压力边界检验；问题二的多目标调度、基线对比、资源能源复算和敏感性检验；问题三的储能优化、区域指标、基准复核和约束检验；问题四的算力--储能--电力联合调度、情景比较、代表窗口检验和独立核验。

\subsection*{C. 结果使用边界}

正文中的数值以对应结果表为准，不从图片估读。通过硬约束检验只能说明当前方案在给定数据、参数和边界下具有可行性，不能单独推出全局最优或对所有场景普遍最优。对于采用数学启发式得到的结果，正文如实说明其求解方法、对照检验和剩余局限。

时间边界统一解释为：0--2399小时为主要运行期，2400--2405小时用于处理已到达任务的尾部执行，2406小时为终端结算时点，不作为普通任务运行小时。成本指标若包含售电收益，应按净结算值解释，不等同于实际用电量为负。

\subsection*{D. AI工具使用详情}

本参赛队在竞赛过程中使用了DeepSeek，主要用于语言润色、代码调试、图表制作与LaTeX排版辅助。AI未用于确定核心模型、公式、参数、原始数据和最终结论。

\subsubsection*{D.1 工具与使用日期}

\begin{table}[H]
    \centering
    \small
    \renewcommand{\arraystretch}{1.05}
    \setlength{\tabcolsep}{3pt}
    \begin{tabularx}{0.96\textwidth}{>{\centering\arraybackslash}p{0.18\textwidth}>{\centering\arraybackslash}p{0.20\textwidth}>{\centering\arraybackslash}X>{\centering\arraybackslash}p{0.22\textwidth}}
        \toprule
        AI工具 & 版本/型号 & 开发机构 & 使用日期 \\
        \midrule
        DeepSeek & DeepSeek-v4-flash-0731 & 深度求索（DeepSeek） & 2026年8月7日至8月10日 \\
        \bottomrule
    \end{tabularx}
\end{table}

本文使用的模型版本为DeepSeek‐v4‐flash‐0731。

\subsubsection*{D.2 使用环节与目的}

DeepSeek用于文字润色、程序调试、图表制作、终稿语言压缩、术语统一、正文引用、结果版本一致性检查和LaTeX排版复核。AI未用于生成或篡改原始数据，未改变既有联合任务搜索的目标函数、约束条件和参数设置；新增结果均由本地程序实际运行并独立复算。

\subsubsection*{D.3 关键交互记录}

\noindent\textbf{交互记录1（DeepSeek辅助润色，使用日期：2026年8月7日至10日）}

提示词要求：只检查由参赛队自行完成的论文文字，不增加新的模型、公式、数据或结论，并给出局部修改建议。AI主要指出长句、重复表述、术语不一致和图表衔接问题；参赛队人工筛选、核对并重新表述。

\noindent\textbf{交互记录2（代码调试与图表辅助，使用日期：2026年8月9日）}

提示词要求：仅检查已有程序的语法、索引、数据类型、文件读取和绘图参数，不修改原始数据、目标函数、约束条件、模型结构、参数设置和最终结论。AI提出变量检查、索引调整、数据类型转换和图表标签优化建议；参赛队仅采纳本地运行可复现且不改变模型含义的建议。

\noindent\textbf{交互记录3（DeepSeek终稿复核，使用日期：2026年8月10日）}

提示词要求：严格依据评审清单核对问题四的结果口径，消除顺序基准、联合方案、图表和正文之间的数据混用，并统一学术语体、文内引用和LaTeX排版。AI协助增加固定任务方案的统一能源定向复算与独立核验入口，修改绘图和论文源文件；所有新指标均由参赛队计算机本地运行生成，参赛队负责最终核对和提交。

\subsubsection*{D.4 采纳与人工修改}

AI建议均经参赛队核对后采纳。所有核心方法、公式、参数、结果和结论由参赛队确定；最终程序、数据、图表和结论由参赛队在本地运行并复核。DeepSeek完整记录见支撑材料中的相应使用报告。

程序中的AI辅助说明对应DeepSeek参与的代码检查与修改环节。AI辅助范围限于程序检查、结果复核、图表制作和论文排版，未改写既有联合任务搜索模型。

\subsection*{E. 源程序清单与附件位置}

本题实际使用的全部24个Python源程序不在论文附录中重复排印，统一收录于独立附件\texttt{CCM2601033附件.zip}的\texttt{源程序}目录。附件同时包含\texttt{pyproject.toml}、\texttt{uv.lock}、处理后的模型输入、关键结果表、模型检验表、论文图形及SHA-256校验清单，赛题提供的原始数据未重复收录。

源程序按问题一至问题四的相对目录组织，共24个\texttt{.py}文件；复核时按附录A所列命令依次运行。论文与附件分开提交，附件命名符合“选题+参赛编号+附件”的要求，压缩包大小控制在20 MB以内。
